\documentclass[10pt,a4paper]{amsart}
\usepackage[margin=19mm]{geometry}
\usepackage{amsmath,amssymb,mathtools}
\usepackage[scr=rsfs,cal=euler]{mathalfa}
\usepackage{xfrac}
\usepackage[unicode,hidelinks,bookmarks=false]{hyperref}
\newcommand{\ie}{i.e.\ }
\newcommand{\cf}{cf.\ }
\newcommand{\resp}{respectively\ }
\newcommand{\Subsection}{\subsection}
\providecommand{\nobreakdash}{\nobreak\hbox{-}}
\setlength{\emergencystretch}{2em}
\multlinegap0pt
\usepackage{bm}
\usepackage{bbm}
\usepackage{dsfont}
\usepackage{xspace}
\usepackage{cancel}
\usepackage{mathtools}
\usepackage{amsthm}
\makeatletter
\@ifundefined{theorem}{\newtheorem{theorem}{Theorem}}{}
\@ifundefined{definition}{\newtheorem{definition}{Definition}}{}
\@ifundefined{lemma}{\newtheorem{lemma}[theorem]{Lemma}}{}
\makeatother
\title[On the number of edges of restricted matchstick graphs]{On the number of edges of restricted matchstick graphs}
\author{Panna Geh\'er, J\'anos Pach, Konrad Swanepoel, G\'eza T\'oth}
\date{Source: arXiv:2506.01589v1 (CC BY 4.0)}
\begin{document}
\maketitle

\noindent\emph{Source and licence.} On the number of edges of restricted matchstick graphs. Version arXiv:2506.01589v1 (CC BY 4.0), distributed under the Creative Commons Attribution 4.0 International licence, as recorded in the arXiv licence metadata for this exact version and shown on the abstract page https://arxiv.org/abs/2506.01589v1. Full rights notice: licenses/arxiv-2506.01589-CC-BY-4.0.txt.\par\smallskip

\noindent\textit{Editorial scope.} Counted unit: the Lemma whose source label is \texttt{monopath} in the article (source0.tex lines 349--351, source label \texttt{monopath}), delivered together with the article's own complete original proof of it (source0.tex lines 402--523, one \texttt{proof} environment of 122 lines). No mathematics is written, repaired or re-derived: statement and proof are the article's own text line for line. The proof is detached from the statement in the source: the article states the result at lines 349--351 and prints its proof at lines 402--523, and the proof environment's own header names the statement, so the association is the article's own. Required context retained uncounted: none. Every definition, display and label of those lines is retained unchanged. Omitted from this package and not redistributed: 1--348, 352--401, 524--629. Nothing inside a retained excerpt is omitted or altered: every retained body is delivered exactly as the article's lines are written, so no omission receipt of any kind accompanies this package. Citations: the retained excerpts carry no citation command at all, so no bibliography entry of the article is transcribed and no reference is redistributed. Reference adaptation: none was needed and none was made; every reference inside the delivered unit resolves inside it, so no unresolved reference or citation is produced. Printed slips kept literal: no printed word, symbol, hypothesis or number of the counted statement or of its proof was altered. Layout and class: the article's own class is \texttt{article} with packages including amssymb, amsmath, amsthm, graphicx, geometry, todonotes, pgfplots, mathrsfs; the wrapper is the lane's pinned \texttt{amsart} editorial wrapper at 10pt on A4, which restates the theorem-like environments the retained text needs and adds the article's own macro definitions that the retained text needs (none), with extra wrapper packages (bm, bbm, dsfont, xspace, cancel, mathtools). The delivered numbering is the unit's own. Assets: no retained span contains a figure environment, a graphics-inclusion command, a PDF-page inclusion, a table or a TikZ picture; no image file is copied into this package, the package declares no asset dependency and no image file is redistributed with it. Licence: version arXiv:2506.01589v1 of this article is distributed under the Creative Commons Attribution 4.0 International licence, as recorded in the arXiv licence metadata for this exact version and shown on the abstract page https://arxiv.org/abs/2506.01589v1. A full rights notice is delivered at licenses/arxiv-2506.01589-CC-BY-4.0.txt. Characters: every retained excerpt is pure ASCII, so the delivered engine, XeLaTeX with its default Latin Modern text font, reports no missing character.
\begin{lemma}\label{monopath}
Any regular edge has an irregular edge at distance at most $16r^2$.
\end{lemma} 

\begin{proof}[Proof of the Lemma \ref{monopath}]
Suppose for a contradiction that there is an edge $\alpha$ 
with the property that its $16r^2$-neighborhood contains only regular edges. Without loss of generality we can assume that $\alpha$ is horizontal; otherwise, we rotate the drawing of $G''$.
\smallskip

A \emph{monotone path} is a path in $G''$ that intersects any vertical line at most once. 
Let $P$ be a monotone path in $G''$ and suppose that $R$ is a 
rhombus-face whose lower left and lower right edges are in $P$. Then $R$ is called a \emph{hat} on $P$. 
\smallskip

We apply the following procedure which builds a monotone path.

\bigskip


\textsc{Algorithm Extend-Path($\alpha$)}

\smallskip

\noindent \textsc{Step $0$.} Let $P_1$ be the path of length $1$ that consists of the edge $\alpha$.
Set $i=2$.


\smallskip

For $i=2, \ldots, $ we do the following.

\noindent \textsc{Step $(i, \, 1)$}. We have $i\ge 2$. 
Suppose that we already have a monotone path $P_{i-1}$.

If there is a \emph{hat} $H$ on it, then replace its two lower edges by its two upper edges in $P_{i-1}$. Let $P_{i-1}$ be the resulting monotone path.
Repeat \textsc{Step $(i, \, 1)$}.

If there is no hat on $P_{i-1}$, then go to \textsc{Step $(i, \, 2)$}.

\smallskip

\noindent \textsc{Step $(i, \, 2)$}. We have a monotone path $P_{i-1}$ of length $i-1$ with no hat on it.
Let $p_0, \, p_1, \, \ldots, \, p_i$ be the vertices of $P_{i-1}$ from left to right and let $F_j$ be the face on the upper side of 
$\epsilon_j=p_{j-1}p_j$. If $F_j$ is not a rhombus, then $\epsilon_j$ is an irregular edge. Let $\epsilon=\epsilon_j$ and {\textbf{Stop.}} \textbf{Return $(\epsilon, \, P_{i-1})$}.

\smallskip
We can assume now that all $F_j$ are rhombi and all of them are different. 
For every $j$, $\epsilon_j$ 
is \emph{rightsided} (resp.\ \emph{leftsided})
if $\epsilon_j$ is the lower \emph{right} (resp.\ lower \emph{left}) side of $F_j$. 
If $\epsilon_1$ is \emph{rightsided}, 
then $P_{i-1}$ can be extended by the lower left edge of $F_1$ to a monotone path $P_i$ 
of length $i$. Let  
$P_{i}$ be this path. Increase $i$ by one and go to \textsc{Step $(i, \, 1)$}.

Now we can assume that $\epsilon_1$ is \emph{leftsided}. 
But then $\epsilon_2$ is also \emph{leftsided}, otherwise 
$F_1$ and $F_2$ would both contain the points slightly above $p_1$, which is impossible since they
internally disjoint. By repeated application of the same argument, we get that all edges 
$\epsilon_1, \, \epsilon_2, \, \ldots, \, \epsilon_{i-1}$ are \emph{leftsided}.

So $\epsilon_{i-1}$ is 
\emph{leftsided},
therefore, $P_{i-1}$ can be extended by the lower right  
edge of $F_{i-1}$ to a monotone path $P_i$ 
of length $i$. Let  
$P_{i}$ be this path. 
Increase $i$  by one and go to \textsc{Step $(i, \, 1)$}.


\bigskip


\begin{definition}
For any monotone path $P$ of 
length $l$, let $c(P)$ denote the \emph{convexity number} of $P$.
It is defined as the number of pairs of edges 
$(\epsilon_1, \epsilon_2)$ of $P$ with the property that 
$\epsilon_1$ is to the left of $\epsilon_2$, and $\epsilon_1$ has smaller slope than $\epsilon_2$.
\end{definition} 

\smallskip

So, for any monotone path $P$ of length $l$, we have
$0\le c(P)\le \binom{l}{2}$. If 
$c(P)=0$, then $P$ is convex from below, if 
$c(P)=\binom{l}{2}$, then 
it is convex from above.

\textsc{Algorithm Extend-Path($\alpha$)}
will terminate after a finite number of steps, since $G''$ contains no monotone path longer than $n$.

Suppose hat the algorithm returns edge $\epsilon$ and path $P_l$ of length $l$.
We distinguish two cases. 

\smallskip

\noindent \textbf{Case 1:} $l\le 4r$.
The algorithm stopped in \textsc{Step $(l+1, \, 2)$}.  

For 
$1\le i\le l+1$, let $s(i)$ be the number of 
times 
\textsc{Step $(i,  \, 1)$} 
was executed. 
Observe that in each step, $c(P)$ decreases by one.
On the other hand, in \textsc{Step $(i,  \, 2)$}, when we construct 
$P_i$ from $P_{i-1}$, we have $c(P_i)\le c(P_{i-1})+i-1.$
Therefore, $$\sum_{i=1}^{l+1}s(i)\le \sum_{i=1}^{l+1}(i-1)<l(l+1)/2.$$ 
Consequently, the total number of steps executed by the algorithm is at most $l(l+1)/2+l\le 16r^2$. 
Thus, 
all edges of $P_{l}$, including $\epsilon$, are  
at distance at most $16r^2$ from $\alpha$, contradicting our assumption that $\epsilon$ is an irregular edge. 

\smallskip


\noindent \textbf{Case 2:} $l\ge 4r$. For $i=4r$, the algorithm  constructed $P_{4r}$, a monotone path of length $4r$.

Again, for $1\le i\le 4r$, let $s(i)$ be the number of times \textsc{Step $(i,  \, 1)$} was executed. By the same argument as in Case 1, we obtain that 
$$\sum_{i=1}^{4r}s(i)\le \sum_{i=1}^{4r}(i-1)<4r(4r-1)/2.$$
Consequently, the total number of steps executed by the algorithm is at most $4r(4r-1)/2+4r<16r^2$. Therefore, 
all edges of $P_{4r}$ are at distance at most $16r^2$ from $\alpha$. By the assumption that there is no irregular edge in the 
$16r^2$-neighborhood of $\alpha$, we conclude that from $\alpha$ to any edge of $P$, there is a sequence of adjacent rhombi, of length at most $16r^2$. 

We also assumed that there is no fat rhombus. Therefore, two edges of a rhombus determine an angle at most $\pi/50r^2$. It follows that the angle between $\alpha$ and any edge of $P_{4r}$ is less than $16r^2\pi/50r^2< \pi/3$. Consequently, the length of the projection of every edge of $P_{4r}$ to the $x$-axis is longer than $1/2$, so the projection of $P_{4r}$ to the $x$-axis is longer than $2r$. This contradicts the fact that $P_{4r}$ lies entirely in $D$, a disk of radius $r$. This concludes the proof of the Lemma. \end{proof}

\end{document}
